\subsubsection{ONT RNA Sample Preparation}

One mL aliquots of actively growing~\jsyn~and~\jsyna~cells were placed in 25 mL of SP4 + 17\% KnockOut serum replacement in 50 mL conical tubes and incubated overnight at 37\textdegree C.
The cells were harvested in mid-late log phase as indicated by a red-orange color of the phenol red pH indicator in the growth media, when the media pH was 7.0.
~\syn~cultures achieved this growth stage in 16 hours and~\syna~cultures in 22 hours.
Growth was arrested in the bulk cultures by placing the tubes on wet ice for 10 minutes.
Keeping all tubes on ice, 1.4 mL of bulk culture was distributed into fifteen 1.8 mL microfuge tubes per culture and centrifuged at 4\textdegree C for five minutes at 10,000$\times$g using a tabletop refrigerated centrifuge to pellet the cells.
The supernatants from ten of the fifteen tubes were aspirated, and in the remaining five tubes the supernatant was used to suspend the pellets.
The cell suspensions from those five tubes were then used to resuspend the pellets in five other tubes, and those more concentrated cell suspensions were used to resuspend the cells in the last five tubes, leaving both the~\syn~and~\syna~cells in five tubes of concentrated cell suspension.
These tubes were centrifuged for five minutes at 4\textdegree C under 10,000$\times$g, the supernatants were aspirated completely, and the remaining pellets were snap frozen using liquid nitrogen, transferred to dry ice, and immediately stored at -80\textdegree C.

Cellular RNA was extracted from~\syn~and~\syna~cell samples using a phenol-chloroform extraction protocol adapted from~\cite{grunberger_next_2019} and the ThermoFisher TRIzol\textsuperscript{TM} Reagent user guide.
Briefly, 1 mL of 4\textdegree C TRIzol\textsuperscript{TM} was added to each tube containing frozen cell pellets, and the liquid was pipetted up and down multiple times to break up the cell pellet.
The tubes were incubated at room temperature for 5 minutes to allow complete disassociation of nucleoprotein complexes.
200 $\mu$L of chloroform was then added and the tubes were mixed by shaking for 30 seconds.
After the tubes incubated at room temperature for 3 minutes, they were centrifuged at 12,000$\times$g for 15 minutes at 4\textdegree C.
The upper RNA-containing aqueous phase was pipetted without collecting any of the white interphase separating the aqueous from the phenol-chloroform lower phase, and placed in fresh 1.8 mL tubes already containing 500 $\mu$L isopropanol.
The tubes were incubated on ice for 10 minutes and then centrifuged at 12,000$\times$g for 10 minutes at 4\textdegree C.
The supernatant was aspirated and discarded off the white gel-like pellets, the tubes were inverted, and the pellets were allowed to air dry for 15 minutes.
The RNA pellets were then dissolved in 20 $\mu$L of RNase-free water containing 0.1 mM EDTA.
RNA concentration was assessed using a Qubit RNA HS fluorometer assay and either an Agilent Bioanalyzer 2100 (\syn) or an Agilent TapeStation 2200 (\syna).
The next step was ribosomal RNA (rRNA) depletion.

The original~\syn~sample was rRNA depleted using an Invitrogen RiboMinus\textsuperscript{TM} Bacteria 2.0 Transcriptome Isolation Kit.
Briefly, 2 $\mu$g of total~\syn~RNA was mixed with 100 $\mu$L of RiboMinus\textsuperscript{TM} 1X hybridization buffer and 6 $\mu$L of RiboMinus\textsuperscript{TM} Pan-Prokaryote Probe Mix, and the solution was brought up to 200 $\mu$L using RNase-free water.
The solution was denatured for 10 minutes at 70\textdegree C and then hybridized for 20 minutes at 37\textdegree C.
The hybridization solution was mixed with RiboMinus\textsuperscript{TM} Magnetic Beads (1 mL of suspension) that had been washed in RNase-free water and resuspended in 400 $\mu$L of RiboMinus\textsuperscript{TM} 1X hybridization buffer, and incubated at 37\textdegree C for 15 minutes.
The sample tube was placed on a RiboMinus\textsuperscript{TM} magnetic stand, and the supernatant (600 $\mu$L) containing the rRNA-depleted RNA was transferred to two new 1.8 mL tubes.
To those tubes, each containing 300 $\mu$L of RNA, premixed Binding Solution Concentrate (400 $\mu$L) and Nucleic Acid Binding Beads (10 $\mu$L) supplied with the kit were added and gently mixed.
One mL of 100\% ethanol was added to each tube, mixed, and incubated at room temperature for 5 minutes.
The tubes were returned to the magnetic stand for 5 minutes, and the supernatant was aspirated and discarded without disturbing the beads.
The beads were washed in 300 $\mu$L of the kit Wash Solution, the tubes were returned to the magnetic stand, and the cleared supernatant was aspirated and discarded without disturbing the bead pellet.
After the beads had dried for 5 minutes, 25 $\mu$L of pre-heated (70\textdegree C) nuclease-free water was added and the tubes were incubated for 1 minute at room temperature to elute the RNA.
The tubes were returned to the magnetic stand and the supernatant was carefully collected into a new microcentrifuge tube.

The second~\syn~sample and the~\syna~sample were rRNA depleted using the NEBNext\textregistered{} rRNA Depletion Kit for bacteria.
One $\mu$g of RNA was brought up to 11 $\mu$L with RNase-free water, and 2 $\mu$L of NEBNext\textregistered{} Bacterial rRNA Depletion Solution and 2 $\mu$L of NEBNext\textregistered{} Probe Hybridization Buffer were added.
The reaction was heated to 95\textdegree C for 2 minutes and cooled to 22\textdegree C at a ramp rate of 0.1\textdegree C/sec.
This was followed immediately by RNase H digestion, where 2 $\mu$L of RNase H Reaction Buffer and 2 $\mu$L of NEBNext Thermostable RNase H were added to the hybridized RNA solution to a total volume of 20 $\mu$L, and incubated at 50\textdegree C for 30 minutes.
Following a bead cleanup with Agencourt RNAClean XP beads (Beckman Coulter\texttrademark, cat \# A63987), the 15 $\mu$L of the rRNA-depleted solution was poly(A) tailed at the 3$'$ end.

The same poly(A) tailing protocol was used for~\syn~and~\syna.
RNA samples were mixed with a premade solution containing 2 $\mu$L of 10X~\eco~Poly(A) Polymerase Reaction Buffer, 2 $\mu$L of ATP, and 1 $\mu$L of~\eco~Poly(A) Polymerase (New England Biolabs) and incubated at 37\textdegree C for 30 minutes.
The reaction was stopped by adding EDTA to a final concentration of 10 mM.
Following another bead cleanup along with the necessary quality control analysis, RNA was normalized to 9 $\mu$L containing 140 ng for~\syn~and 50 ng for~\syna~RNA samples.

\subsubsection{ONT RNA Sample Quality Control Analysis}

An Agilent Bioanalyzer 2100 assay was used for assessing RNA quality and concentration of the first~\syn~sample both before and after rRNA depletion.
For the second~\syn~and the~\syna~samples, both a Qubit RNA HS fluorometer assay and an Agilent TapeStation 2200 were used for assessing RNA concentration post-extraction from minimal cells.
Initial Qubit values post RNA extraction were 3000 ng/$\mu$L for the second~\syn~and 1890 ng/$\mu$L for~\syna.
The Agilent TapeStation 2200 was used throughout the extraction and rRNA depletion steps to assess the quality of the RNA and depletion, giving an estimate of depletion through the absence of the 16S and 23S peaks (\sdqc).
A Qubit RNA HS fluorometer was used to detect the final quantitation of RNA.
Post poly-adenylation of RNA transcripts, the extent of DNA and buffer contamination in each of the RNA samples was tested by performing standard spectroscopic measurements (NanoDrop One) and using the Qubit 1$\times$ dsDNA HS assay kit (ThermoFisher Scientific).
Input RNA samples for Oxford Nanopore Direct RNA Sequencing library preparation were finally quantified using the Qubit RNA HS assay kit, and RNA concentration values were used to normalize samples to the desired input volume of 9 $\mu$L containing 50--500 ng of input RNA for library preparation.

\subsubsection{ONT Library Preparation and Direct Sequencing}

After normalization, libraries for Oxford Nanopore sequencing were prepared using the ONT Direct RNA Sequencing Kit (SQK-RNA002) with slight modifications.
Specifically, 0.5 $\mu$L of the RT adapter enzyme (reduced from 1 $\mu$L) and 4 $\mu$L of the RMX RNA adapter (reduced from 6 $\mu$L) were used during library preparation for~\syn~and~\syna.
All samples were sequenced using R9.4 flow cells on the ONT GridION platform, with live base-calling enabled through the recommended MinKNOW (v23.07.5) scripts to generate fast5 files.
For~\syn, 819{,}435 reads were generated over two runs, and for~\syna~734~k reads over a single run.
In all run reports, base-called bases were 99.99--100\% confident.

\subsubsection{Base-calling of Raw ONT Reads}

Results of the Nanopore reads were collected as individual fast5 files, which were combined into multi-read files for analysis, and current signals were converted to canonical bases by base-calling.
For the first~\syn~sequencing, base-calling was performed with Guppy (v6.1.3; Oxford Nanopore Technologies, Ltd., 2022) using the RNA trim strategy and with calibration detection and filter enabled, on NVIDIA A5000 GPUs.
For the second~\syn~replicate and the~\syna~experiments, Guppy (v7.0.9, GUI-based for the MinKNOW system) was used for base-calling with equivalent parameters.

\subsubsection{Mapping ONT Reads to the Reference Genomes}

Reads were retrieved from the NCBI SRA with the SRA Toolkit:~\syn~run 1 (SRR36199726, 234{,}477 reads),~\syn~run 2 (SRR36199725, 584{,}958 reads), and~\syna~(SRR36199724, 734~k reads).
Direct-RNA reads deposited in the native 3$'\rightarrow$5$'$ pore order (the~\syn~run 2 and the~\syna~run) were restored to 5$'\rightarrow$3$'$ sense orientation by reversing the base order (\texttt{seqkit} v2.1.0, reverse only and not reverse-complement, since direct-RNA reads are single-stranded), whereas the~\syn~run 1 reads were already in sense orientation.
Orientation was verified empirically: in the corrected orientation reads mapped sense to the annotated genes (over 97\% of gene-overlapping reads), whereas the opposite orientation mapped almost nothing.
Reads were aligned to the respective reference genome (\jsyn, NCBI accession CP002027.1;~\jsyna, CP016816.2) with \texttt{minimap2} v2.30~\cite{li_minimap2_2018} using the non-spliced long-read preset \texttt{-ax map-ont -p 0.99 --MD}, and sorted and indexed with \texttt{samtools} v1.22.1~\cite{li_sequence_2009}.
The \texttt{map-ont} preset was used in place of the splice-aware preset because~\mmy~and its reduced derivative are intron-less, so N-skip (spliced) alignments would bridge tens to hundreds of kilobases of unrelated reference.
Strandedness of the direct-RNA reads is preserved through the BAM flag bits, and per-strand sequencing-depth bedGraph tracks were computed with \texttt{samtools depth}, with forward-mapped reads defining the plus-strand transcripts.

For~\syn, run 1 yielded 207{,}480 primary mapped reads (88.5\%; mean length 661~nt, mean Phred quality Q17.5) and run 2 yielded 456{,}755 primary mapped reads (78.1\%; mean length 370~nt, Q31.7, 98.4\% above Q20), the higher quality of run 2 reflecting its newer base-caller.
For~\syna, 559~k of 734~k reads (76.2\%) mapped as primary alignments (mean length 383~nt, range 49--2{,}858~nt, mean quality Q31.2, 98.2\% above Q20); the 175~k unmapped reads (23.9\%) were of substantially lower quality (mean Q20.7, 61.0\% shorter than 300~nt), consistent with the expected tail of truncated ONT direct-RNA molecules rather than reference mismatch.
Secondary alignments accounted for 3--4\% of primary mapped reads in every run and clustered at the ribosomal RNA and tRNA loci of each genome (for~\syna, $\sim$55{,}460 and $\sim$343{,}267~bp, the two rRNA operons).

\subsubsection{Between-run Reproducibility and TPM Quantification}

Per-gene sense transcripts-per-million (TPM) were computed from the strand-specific depth tracks using the depth-based definition applied to the~\syn~Illumina and PacBio data.
The two~\syn~runs agreed only moderately (Pearson $r = 0.82$ on $\log_{10}$ TPM, with 36\% of genes differing more than two-fold), the largest discrepancies being residual rRNA transcripts, which the RiboMinus\textsuperscript{TM}-depleted run 1 retained far more abundantly than the NEBNext\textregistered{}-depleted run 2.
Because the two runs used different depletion chemistries and differ substantially, they were treated as independent runs rather than averaged into a single quantitative profile, and the pooled alignment was retained only as a genome-browser coverage track.
Each~\syn~run reproduced the~\syn~Illumina TPM profile with comparable fidelity (Pearson $r = 0.61$ and $0.60$ on $\log_{10}$ TPM for runs 1 and 2), confirming that the direct-RNA data recover the same transcriptome as the short-read and PacBio platforms.
The~\syna~ONT TPM profile was quantified in the same way and is reported alongside the~\syna~Illumina track (see the~\syna~short-read sequencing methods).
